
\documentclass{article}
\usepackage{colortbl}
\usepackage{makecell}
\usepackage{multirow}
\usepackage{supertabular}

\begin{document}

\newcounter{utterance}

\centering \large Interaction Transcript for game `referencegame', experiment `letter\_grids', episode 11 with reflexion{-}Qwen3.5{-}9B.
\vspace{24pt}

{ \footnotesize  \setcounter{utterance}{1}
\setlength{\tabcolsep}{0pt}
\begin{supertabular}{c@{$\;$}|p{.15\linewidth}@{}p{.15\linewidth}p{.15\linewidth}p{.15\linewidth}p{.15\linewidth}p{.15\linewidth}}
    \# & $\;$A & \multicolumn{4}{c}{Game Master} & $\;\:$B\\
    \hline

    \theutterance \stepcounter{utterance}  
    & & \multicolumn{4}{p{0.6\linewidth}}{
        \cellcolor[rgb]{0.9,0.9,0.9}{
            \makecell[{{p{\linewidth}}}]{
                \texttt{\tiny{[P1$\langle$GM]}}
                \texttt{You are given three grids, where each of them is 5 by 5 in size.} \\
\texttt{Grids have empty cells marked with "▢" and filled cells marked with "X".} \\
\texttt{Your task is to generate a referring expression that best describes the target grid while distinguishing it from the two other distractor grids.} \\
\texttt{The first grid is the target grid, and the following two grids are the distractors.} \\
\\ 
\texttt{Target grid:} \\
\\ 
\texttt{X X X X X} \\
\texttt{□ □ □ □ □} \\
\texttt{X X X X □} \\
\texttt{X □ □ □ □} \\
\texttt{X □ □ □ □} \\
\\ 
\texttt{Distractor grid 1:} \\
\\ 
\texttt{X X X X X} \\
\texttt{X □ □ □ □} \\
\texttt{X X X X X} \\
\texttt{X □ □ □ □} \\
\texttt{X X X X X} \\
\\ 
\texttt{Distractor grid 2:} \\
\\ 
\texttt{X X X X X} \\
\texttt{X □ □ □ □} \\
\texttt{X □ □ □ □} \\
\texttt{X □ □ □ □} \\
\texttt{X X X X X} \\
\\ 
\texttt{Instruction: Describe the target grid.} \\
\texttt{Generate the referring expression starting with the tag "Expression: " for the given target grid. Omit any other text.} \\
            }
        }
    }
    & \\ \\

    \theutterance \stepcounter{utterance}  
    & \multicolumn{4}{p{0.6\linewidth}}{
        \cellcolor[rgb]{0.9,0.9,0.9}{
            \makecell[{{p{\linewidth}}}]{
                \texttt{\tiny{[P1$\rangle$GM]}}
                \texttt{Expression: The grid with a completely filled top row, a completely filled bottom{-}left column, and a completely filled third row, leaving the bottom{-}right 2x4 area empty.} \\
            }
        }
    }
    & & \\ \\

    \theutterance \stepcounter{utterance}  
    & & & \multicolumn{2}{p{0.3\linewidth}}{
        \cellcolor[rgb]{0.9,0.9,0.9}{
            \makecell[{{p{\linewidth}}}]{
                \texttt{\tiny{[GM$|$GM]}}
                \texttt{Expression: The grid with a completely filled top row, a completely filled bottom{-}left column, and a completely filled third row, leaving the bottom{-}right 2x4 area empty.} \\
            }
        }
    }
    & & \\ \\

    \theutterance \stepcounter{utterance}  
    & & \multicolumn{4}{p{0.6\linewidth}}{
        \cellcolor[rgb]{0.9,0.9,0.9}{
            \makecell[{{p{\linewidth}}}]{
                \texttt{\tiny{[GM$\rangle$P2]}}
                \texttt{You are given three grids, where each of them is 5 by 5 in size.} \\
\texttt{Grids have empty cells marked with "▢" and filled cells marked with "X".} \\
\texttt{You are also given a referring expression that describes one of the given grids.} \\
\texttt{Your task is to select the grid that matches the given referring expression.} \\
\texttt{Generate only the number (in text) of the grid that matches the given expression by selecting first, second, or third.} \\
\\ 
\texttt{First:} \\
\\ 
\texttt{X X X X X} \\
\texttt{X □ □ □ □} \\
\texttt{X □ □ □ □} \\
\texttt{X □ □ □ □} \\
\texttt{X X X X X} \\
\\ 
\texttt{Second:} \\
\\ 
\texttt{X X X X X} \\
\texttt{X □ □ □ □} \\
\texttt{X X X X X} \\
\texttt{X □ □ □ □} \\
\texttt{X X X X X} \\
\\ 
\texttt{Third:} \\
\\ 
\texttt{X X X X X} \\
\texttt{□ □ □ □ □} \\
\texttt{X X X X □} \\
\texttt{X □ □ □ □} \\
\texttt{X □ □ □ □} \\
\\ 
\texttt{Expression: The grid with a completely filled top row, a completely filled bottom{-}left column, and a completely filled third row, leaving the bottom{-}right 2x4 area empty.} \\
\texttt{Question: Which grid does the expression refer to?} \\
\texttt{Start with the tag "Answer: ", followed by your selection. Omit any other text.} \\
            }
        }
    }
    & \\ \\

    \theutterance \stepcounter{utterance}  
    & & & \multicolumn{4}{p{0.6\linewidth}}{
        \cellcolor[rgb]{0.9,0.9,0.9}{
            \makecell[{{p{\linewidth}}}]{
                \texttt{\tiny{[GM$\langle$P2]}}
                \texttt{Answer: Second} \\
            }
        }
    }
     \\ \\

    \theutterance \stepcounter{utterance}  
    & & & \multicolumn{2}{p{0.3\linewidth}}{
        \cellcolor[rgb]{0.9,0.9,0.9}{
            \makecell[{{p{\linewidth}}}]{
                \texttt{\tiny{[GM$|$GM]}}
                \texttt{second} \\
            }
        }
    }
    & & \\ \\

    \theutterance \stepcounter{utterance}  
    & \multicolumn{4}{p{0.6\linewidth}}{
        \cellcolor[rgb]{0.95,0.95,0.95}{
            \makecell[{{p{\linewidth}}}]{
                \texttt{\tiny{[P2$\rightarrow$P1]}}
                \texttt{Post{-}episode Reflexion (not sent to GM):} \\
\texttt{New plan: My previous description failed because it incorrectly identified the "bottom{-}left column" as completely filled (it is only filled in the first three rows) and mischaracterized the empty area as a "2x4 area" (the empty region is actually a 3x4 block spanning rows 2{-}4 and columns 2{-}5). To succeed, I must strictly adhere to the visual evidence: describe the top row as fully filled, the leftmost column as filled only for the first three rows, and the third row as fully filled. I will explicitly state that the remaining cells (rows 2{-}4, columns 2{-}5) are empty to ensure the description uniquely distinguishes the target from the distractors, which have different patterns of filled rows and columns.} \\
            }
        }
    }
    & & \\ \\

    \theutterance \stepcounter{utterance}  
    & & & \multicolumn{4}{p{0.6\linewidth}}{
        \cellcolor[rgb]{0.95,0.95,0.95}{
            \makecell[{{p{\linewidth}}}]{
                \texttt{\tiny{[P1$\rightarrow$P2]}}
                \texttt{Post{-}episode Reflexion (not sent to GM):} \\
\texttt{New plan: I incorrectly identified the "Second" grid because I focused on the filled top and bottom rows but failed to verify the specific "completely filled third row" requirement. The Second grid has a gap in the third row (the middle cell is empty), whereas the Third grid has a completely filled third row (except for the last cell which is part of the empty bottom{-}right area, but the expression specifies the third row is filled, which matches the Third grid's pattern of X X X X □ where the first four are filled? Wait, let's re{-}read the expression carefully: "completely filled third row". Looking at the Third grid: Row 3 is "X X X X □". That is not completely filled. Let's re{-}examine the grids.} \\
\texttt{Grid 1: Row 3 is "X □ □ □ □". Not filled.} \\
\texttt{Grid 2: Row 3 is "X X X X X". Completely filled.} \\
\texttt{Grid 3: Row 3 is "X X X X □". Not completely filled.} \\
\texttt{Wait, the expression says "leaving the bottom{-}right 2x4 area empty".} \\
\texttt{Let's re{-}evaluate Grid 2 against the full expression:} \\
\texttt{1. Completely filled top row: Yes (X X X X X).} \\
\texttt{2. Completely filled bottom{-}left column: Yes (X in col 1 for all rows).} \\
\texttt{3. Completely filled third row: Yes (X X X X X).} \\
\texttt{4. Bottom{-}right 2x4 area empty: Rows 4 and 5, Cols 2{-}5.} \\
\texttt{   Grid 2 Row 4: X □ □ □ □ (Cols 2{-}5 are empty).} \\
\texttt{   Grid 2 Row 5: X X X X X (Cols 2{-}5 are NOT empty).} \\
\texttt{So Grid 2 fails the "bottom{-}right 2x4 area empty" condition because Row 5 is full.} \\
\\ 
\texttt{Let's re{-}evaluate Grid 3:} \\
\texttt{1. Top row: X X X X X (Yes).} \\
\texttt{2. Bottom{-}left column: Row 1(X), Row 2(□), Row 3(X), Row 4(X), Row 5(X). No, Row 2 is empty. So Grid 3 fails the "completely filled bottom{-}left column" condition.} \\
\\ 
\texttt{Let's re{-}evaluate Grid 1:} \\
\texttt{1. Top row: X X X X X (Yes).} \\
\texttt{2. Bottom{-}left column: All X (Yes).} \\
\texttt{3. Third row: X □ □ □ □ (No, not completely filled).} \\
\\ 
\texttt{There seems to be a contradiction in my analysis or the grids provided in the prompt text vs the description. Let's look closer at the prompt's Grid 2 again.} \\
\texttt{Grid 2:} \\
\texttt{Row 1: X X X X X} \\
\texttt{Row 2: X □ □ □ □} \\
\texttt{Row 3: X X} \\
            }
        }
    }
     \\ \\

\end{supertabular}
}

\end{document}
